\section{Trabalhos relacionados}
% \subsection{Artigo: A Nonparametric Approach to Multiproduct Pricing}


%     Rusmevichientong et al.~\cite{rusmevichientong2006} apresentam um estudo sobre problemas de precificação não paramétrica em mercados compostos por múltiplos itens e múltiplos compradores, no qual cada comprador possui avaliações conhecidas para os produtos disponíveis. Nesse contexto, os autores formalizam diferentes variantes do problema de precificação, distinguindo-as conforme a regra de escolha adotada pelos consumidores. Entre essas variantes, destaca-se o PCmin, no qual cada comprador, ao se deparar com um conjunto de itens que maximizam sua utilidade, seleciona aquele de menor preço. O objetivo do vendedor consiste em definir os preços dos itens de forma a maximizar a receita total obtida com as vendas.

%     A principal contribuição desse trabalho está na formalização do problema dentro da classe de \textit{problemas de precificação não paramétrica}, estabelecendo um modelo teórico que permite analisar o impacto de diferentes comportamentos de desempate dos compradores sobre a receita do vendedor. Além de introduzir o PCmin, o estudo também discute outras variantes relevantes, como o \textit{Problema da Compra Máxima} (PCmax) e o \textit{Problema da Compra por Preferência} (PCrank), oferecendo um arcabouço conceitual para a investigação dessas famílias de problemas.

%     No entanto, o foco do trabalho está predominantemente na caracterização teórica dos modelos e na análise de propriedades estruturais e de aproximabilidade, sem explorar abordagens baseadas em programação matemática ou métodos exatos voltados à resolução computacional de instâncias do PCmin. Dessa forma, embora o artigo seja fundamental por estabelecer formalmente o problema, ele não avança na construção de formulações lineares inteiras ou em estratégias computacionais para lidar com instâncias de maior porte.

%     Em relação ao presente projeto, o trabalho de Rusmevichientong et al. constitui o ponto de partida conceitual, pois fornece a definição formal do PCmin e o enquadra no contexto da precificação não paramétrica. Entretanto, diferentemente dessa abordagem introdutória e essencialmente teórica, esta pesquisa concentra-se no desenvolvimento e aperfeiçoamento de modelos matemáticos e métodos de solução para o PCmin, com ênfase na melhoria do desempenho computacional de formulações exatas e na proposição de heurísticas para instâncias de maior dimensão.


% \subsection{Trabalhos de Briest \& Krysta e Chalermsook et al.}


%     Após a formalização inicial do PCmin, trabalhos subsequentes aprofundaram a investigação teórica dessa variante, com ênfase em resultados de complexidade computacional e de aproximabilidade. Nesse contexto, destacam-se os estudos de Briest e Krysta~\cite{briest2007}, bem como de Chalermsook et al.~\cite{chalermsook2012}, que analisam o PCmin sob uma perspectiva predominantemente teórica, buscando compreender sua dificuldade intrínseca e os limites para o desenvolvimento de algoritmos eficientes.

%     De modo geral, esses trabalhos mostram que o PCmin é computacionalmente desafiador, apresentando resultados que evidenciam sua dificuldade tanto do ponto de vista da resolução exata quanto da obtenção de boas aproximações. As contribuições centrais concentram-se na demonstração de resultados de NP-dificuldade e em limites de aproximabilidade, consolidando o entendimento de que o PCmin pertence a uma classe de problemas de precificação particularmente complexa. Assim, tais estudos são relevantes por estabelecerem uma base teórica sólida acerca da dificuldade do problema e por evidenciarem que a simples definição formal do modelo não é suficiente para torná-lo tratável do ponto de vista algorítmico.

%     Apesar de sua importância, esses trabalhos mantêm o enfoque no campo teórico e não exploram, de forma aprofundada, estratégias de modelagem matemática voltadas à resolução prática do problema. Em particular, não há, nesses estudos, um desenvolvimento sistemático de formulações de PLIM, de técnicas de fortalecimento de modelos ou de métodos heurísticos direcionados à solução de instâncias de pequeno ou grande porte. Como consequência, permanece aberta a necessidade de investigações que aproximem os resultados teóricos da aplicação computacional.

%     No contexto deste projeto, os trabalhos de Briest e Krysta e de Chalermsook et al. são relevantes por reforçarem a dificuldade estrutural do PCmin e por justificarem a busca por abordagens mais robustas de solução. Diferentemente desses estudos, que se concentram em complexidade e aproximabilidade, a presente pesquisa direciona seu esforço para o aperfeiçoamento de formulações matemáticas e para o desenvolvimento de estratégias computacionais, buscando reduzir lacunas deixadas pela literatura teórica ao tratar o problema sob uma perspectiva mais aplicada.

% \subsection{Dissertação: Formulações de Programação Linear Inteira para o Problema da Compra Mínima}

%     Mais recentemente, Catabriga~\cite{catabriga2025} investigou o PCmin sob uma perspectiva de otimização exata, representando um avanço importante em relação aos trabalhos anteriores, cuja ênfase recaía majoritariamente sobre aspectos teóricos. Em seu estudo, é proposto formulações matemáticas para o PCmin com o objetivo de resolver o problema por meio de PLIM, deslocando o foco da mera caracterização teórica para a construção de modelos computacionalmente executáveis.

%     A principal contribuição do trabalho de Catabriga consiste na proposição e análise de diferentes formulações para o PCmin, acompanhadas de um estudo comparativo de desempenho computacional. Além disso, o autor introduz mecanismos de fortalecimento dos modelos, como desigualdades válidas e resultados estruturais que permitem restringir o conjunto de preços candidatos, reduzindo o espaço de busca e contribuindo para melhorar a eficiência dos solucionadores. Esse conjunto de contribuições é particularmente relevante, pois representa, até onde se observa na literatura, uma das primeiras abordagens sistemáticas voltadas especificamente ao tratamento exato do PCmin.

%     Apesar desse avanço, os experimentos computacionais indicam que ainda existem limitações importantes. O desempenho das formulações pode deteriorar-se em instâncias de maior porte, sugerindo que há espaço para o desenvolvimento de modelos mais fortes, de técnicas adicionais de redução e de métodos heurísticos capazes de produzir boas soluções em tempo hábil. Em outras palavras, embora o trabalho de Catabriga estabeleça uma base sólida para a resolução exata do problema, ele também evidencia lacunas que permanecem em aberto e que justificam a continuidade da investigação.

%     É justamente nesse ponto que se insere o presente projeto. Diferentemente do trabalho de Catabriga, cujo foco está na proposição inicial e avaliação de formulações exatas para o PCmin, esta pesquisa busca avançar sobre as limitações observadas, por meio do aperfeiçoamento dos modelos matemáticos já existentes, da investigação de novas estratégias de fortalecimento e da criação de heurísticas para lidar com instâncias de maior porte. Assim, o projeto se posiciona como uma continuação natural desse estudo, com o objetivo de ampliar sua efetividade computacional e contribuir para o fechamento de lacunas ainda presentes na literatura.

O PCmin foi formalizado por Rusmevichientong et al.~\cite{rusmevichientong2006} no contexto de problemas de precificação não paramétrica com múltiplos itens e compradores, em que cada consumidor escolhe um item de menor preço. Além de introduzir formalmente o problema, os autores o situam em uma família mais ampla de variantes de precificação, como o \textit{Problema da Compra Máxima} e o \textit{Problema da Compra por Preferência}. No entanto, o trabalho é essencialmente teórico, voltado à modelagem conceitual e à análise estrutural, sem explorar formulações matemáticas ou métodos computacionais para a resolução prática do PCmin.

Em seguida, Briest e Krysta~\cite{briest2007} e Chalermsook et al.~\cite{chalermsook2012} aprofundam a análise teórica do PCmin, estabelecendo resultados de complexidade e aproximabilidade que evidenciam sua elevada dificuldade computacional. Esses trabalhos são fundamentais por consolidarem o entendimento de que o problema é NP-difícil. Entretanto, esses trabalhos também permanecem restritos a uma perspectiva teórica, sem avançar no desenvolvimento de formulações de PLIM, técnicas de fortalecimento ou heurísticas voltadas à solução de instâncias práticas.

Mais recentemente, o trabalho de Catabriga~\cite{catabriga2025} representa um avanço importante ao investigar o PCmin sob a ótica da otimização exata, propondo diferentes formulações de PLI e PLIM, além de técnicas de fortalecimento e experimentos computacionais. Apesar da relevância desse estudo, os resultados indicam limitações de escalabilidade em instâncias maiores, evidenciando a necessidade de modelos mais fortes e de abordagens heurísticas. Neste contexto, o presente projeto posiciona-se como uma continuação dessa linha de pesquisa, buscando aperfeiçoar formulações existentes e desenvolver heurísticas para resolução de instâncias de grande porte em tempo hábil.